\chapter{Panneau de réglages en direct}
Changer un réglage d'émulation, puis revenir au morceau pour l'écouter, ne permet pas de comparer : l'oreille oublie
vite. Le panneau en direct règle ce problème. Un appui long sur l'oscilloscope, ou le bouton des réglages, l'ouvre
par-dessus la visualisation pendant que la musique continue.

Il ne montre que ce qui concerne le morceau en cours : les réglages de la puce jouée, ceux de l'oscilloscope et les
réglages généraux. Vous pouvez donc passer de reSIDfp à reSID, ou d'un filtre 6581 à un autre, et entendre la
différence dans la seconde. Les rares réglages qui demandent de recharger le morceau sont signalés par la mention
« applied to the next track ».
\capturedouble{settings-sid}{settings-general}{Page SID}{page Général}
